As described, sharpening in the sense of \cref{def:sharpening} is a
purely computational problem, which makes it difficult to evaluate the
quality and optimality of
self-improvement algorithms. To address this, we introduce a novel
statistical
framework for sharpening, inspired by the success of oracle complexity in optimization
\citep{nemirovski1983problem,traub1988information,raginsky2011information,agarwal2012information}
and statistical query complexity in computational learning theory \citep{blum1994weakly,kearns1998efficient,feldman2012complete,feldman2017general}.


\begin{definition}[Sample-and-evaluate framework]\label{def:oracle-model}
  In the \textbf{\framework} framework, the algorithm designer does
  not have explicit access to the base model $\piref$. Instead, they
access $\piref$ only through \emph{sample-and-evaluate queries}: The learner is allowed to sample $n$ prompts $x \sim
\cdist$. For each prompt $x$, they can sample $N$
responses $y_1,y_2,\dots y_N \sim \piref(\cdot \mid x)$ and observe
the likelihood $\piref(y_i\mid{}x)$ for each such response. The
efficiency, or \emph{sample complexity}, of the algorithm is measured through the total number of
sample-and-evaluate queries
$m\ldef{}n\cdot{}N$.
\end{definition}
This framework can be seen to capture algorithms like \bestofnalg and
\rlhfalg (implemented with DPO\dfdelete{ or PPO}), which only
access the base model $\piref$ through i) sampling responses via
$y\sim\piref(\cdot\mid{}x)$ \textbf{(generation)}, and ii) evaluating the likelihood
$\piref(y\mid{}x)$ \textbf{(verification)} for these responses. We view the sample complexity
$m=n\cdot{}N$ as a natural statistical abstraction for the
computational complexity of self-improvement (a clear parallel to oracle complexity for optimization algorithms), one which is amenable to
information-theoretic lower bounds.\footnote{Concretely, the sample
  complexity $m=n\cdot{}N$ is a lower bound on the running time
  of any algorithm that operates in the \framework framework.} We will
aim to show that, under appropriate assumptions, \bestofnalg and
\rlhfalg can learn an $(\eps,\delta)$-sharpened model with 
sample complexity
\[
  m = \poly(\eps^{-1},\delta^{-1},\Cprob)
\]
where $\Cprob$ is a potentially problem-dependent constant.





\subsection{Fundamental Limits}
\label{sec:lower}

Before diving into our analysis of \bestofnalg and \rlhfalg in the
\framework framework, let us take a brief detour to give a sense for
how sample complexity guarantees for sharpening should scale. To this end, we will prove a lower bound or fundamental limit on
the sample complexity of any algorithm in the \framework framework.

Intuitively, the performance of any sharpening algorithm based on
sampling should depend on how well the base model $\piref$
covers the arg-max response $\ystar(x)$. To capture this, we define
the following \emph{coverage coefficient}:\footnote{This quantity can
  be interpreted as a special case of the $L_1$-concentrability
  coefficient \citep{farahmand2010error,xie2020q,zanette2021provable,amortila2024scalable}
  studied in the theory of offline reinforcement learning.}\loose
\begin{align}
  \label{eq:cstar}
  \Cstar = \En_{x\sim\cdist}\brk*{\frac{1}{\piref(\Ystar(x)\mid{}x)}}.
\end{align}
More generally, for a model $\pi$, we define
$\Ypi(x)=\argmax_{y\in\cY}\pi(y\mid{}x)$ and 
$  \Cstar(\pi) =
\En_{x\sim\cdist}\brk*{\frac{1}{\pi(\Ypi(x)\mid{}x)}}$. 

Our main lower
bound shows that for worst-case choice of $\Pi$, the coverage coefficient acts as a lower bound on the sample
complexity of any sharpening algorithm.
\begin{theorem}[Lower bound for sharpening]
  \label{thm:lower}
  Fix an integer $d \ge 1$ and parameters $\epsilon \in (0,1)$ and $C
  \ge 1$. There exists a class of models $\Pi$ such that (i) $\log |\Pi|
  \asymp d (1+\log(C \epsilon^{-1}))$, (ii) $\sup_{\pi \in \Pi}
  \Cstar(\pi) \lesssim C$, and (iii) $\Ypi(x)$ is a singleton for all
  $\pi\in\Pi$, $x\in\cX$. Any sharpening algorithm $\pihat$ that
  achieves $\En\brk*{\Pr_{x \sim \cdist}[\pihat(\Ypi[\piref](x)\mid{}x) > 1/2]}  \ge 1
  - \epsilon$ for all $\piref \in \Pi$ must collect a total number of
  samples $m = n\cdot{}N$ at least\loose
  \begin{align}
    m \gtrsim
    \frac{ C\log |\Pi|}{\epsilon^{2}\cdot{}(1+\log(C \epsilon^{-1}))}.
  \end{align}  
\end{theorem}
This result shows that the complexity of any $(\eps,1/2-\delta)$-sharpening algorithm (for $\delta>0$) in the
\framework framework must depend polynomially on the coverage
coefficient $\Cstar$, as well as the accuracy\arxiv{ parameter} $\eps$. The
lower bound also depends on the expressivity of $\piref$,
as captured by the model class complexity term $\log\abs{\Pi}$. We will show in \abedit{the sequel} that it is possible to match this lower bound. \dfedit{Note
  that this result also implies a lower bound for the general
    sharpening problem (i.e., general $\rself$), since \mlsharp is a
  special case.}

\begin{remark}[Relaxed notions of sharpening and coverage]
The notion of coverage in \cref{eq:cstar} is somewhat stringent, since
it requires that $\piref$ place large mass on $\Ystar(x)$ on average.
\iclr{In
\cref{sec:proof_preliminaries}, we introduce a more general and permissive
notion of \emph{approximate sharpening}
(\cref{def:sharpening_general}) which
leads to weaker coverage requirements, and use this to
give generalized versions of our main results.\loose
}
\arxiv{In
\cref{sec:proof_preliminaries}, we introduce a more general and permissive
notion of approximate sharpening (\cref{def:sharpening_general}),
which allows the model to sharpen toward approximate arg-max responses
(in the sense that
$\log\piref(y\mid{}x)\geq{}(1-\gamma)\max_{y\in\cY}\log\piref(y\mid{}x)$
for an approximation parameter $\gamma>0$). This notion of sharpening
leads to significantly weaker coverage requirements, and we state
generalized versions of all our main results which accommodate this in
the appendix.
}
\end{remark}

\dfedit{We close this section by noting that numerous
  recent works---focusing on inference-time computation---show that
standard language models exhibit favorable coverage \abedit{with
  respect to desirable responses}
\citep{brown2024large,snell2024scaling,wu2024empirical}. \akedit{We
  replicate these findings in our experimental setup
  in~\Cref{app:experiments}.} These works suggest that, despite the
exponentially large response space, the coverage coefficient $\Cstar$
may be small in standard language modeling tasks.}



    
